\section{Derivations and the Measurement Clock}

\subsection{State and update indexing}
State $\theta_t$ has $t$ completed updates. The raw gradient $g_t$ is evaluated at that state and is available before the outgoing increment $u_t=\theta_{t+1}-\theta_t$. Paired gradient quantities $\rho_t,r_t,\eta\widehat\kappa_t$ compare states $t-1$ and $t$. Representation motion $\delta_t^\ell=\operatorname{vec}(A_t^\ell-A_{t-1}^\ell)$ is incoming and is displayed at its endpoint $t$. In a fitted window $a{:}b$, the raw-gradient comparison uses $g_{a-1},\ldots,g_{b-1}$, the vectors used for those incoming updates. Earlier source records that number a gradient by the update it supplies are shifted by one when displayed; their numerical values are unchanged.

\subsection{Displacement and directional curvature}
For constant-rate plain SGD, summation gives $\theta_T-\theta_0=-\eta\sum_{t=0}^{T-1}g_t$. Expanding its squared norm gives the main all-pair identity. The triangle inequality gives $0\leq R_T\leq1$ when the total path length is nonzero. For another optimizer, the same geometric identity holds for the actual increments $u_t$, but replacing them by unweighted raw gradients generally does not recover parameter displacement.

For any optimizer, let $u_{t-1}=\theta_t-\theta_{t-1}\ne0$ denote the actual parameter change, and let $f_t=\nabla\mathcal L(\theta_t)$ be the full gradient. For a twice continuously differentiable objective, the fundamental theorem gives
\begin{equation}
 f_t-f_{t-1}=\int_0^1H(\theta_{t-1}+a u_{t-1})u_{t-1}\,da.
\end{equation}
Multiplying by $u_{t-1}^{\top}/\|u_{t-1}\|^2$ proves the step-direction curvature identity in the main text. For sampled gradients $g_t=f_t+\xi_t$, it becomes
\begin{equation}
 \widehat\kappa_t^{\mathrm{step}}
 =\int_0^1\frac{u_{t-1}^{\top}H(\theta_{t-1}+a u_{t-1})u_{t-1}}{\|u_{t-1}\|^2}\,da
 +\frac{u_{t-1}^{\top}(\xi_t-\xi_{t-1})}{\|u_{t-1}\|^2}.
 \label{eq:s-general-curvature}
\end{equation}
The identity does not depend on the choice of optimizer or require the step to be parallel to the gradient. Adam, momentum, clipping, and weight decay may change $u_{t-1}$, but the same formula applies to the actual resulting displacement. For a group-specific learning rate, use the complete applied increment rather than multiplying the entire gradient by one nominal rate.

For full gradients and an $L_H$-Lipschitz Hessian in operator norm along the segment,
\begin{align}
 \left|\widehat\kappa_t^{\mathrm{step}}-
 \frac{u_{t-1}^{\top}H(\theta_{t-1})u_{t-1}}{\|u_{t-1}\|^2}\right|
 &\leq\int_0^1\|H(\theta_{t-1}+a u_{t-1})-H(\theta_{t-1})\|_{\rm op}\,da\notag\\
 &\leq\int_0^1L_Ha\|u_{t-1}\|\,da
 =\tfrac12L_H\|u_{t-1}\|.
\end{align}
This specifies when the measured step-average approximates curvature at its starting point. With sampled gradients, the same bound applies after subtracting the explicit noise term in Eq.~\eqref{eq:s-general-curvature}. Across an activation boundary the smooth assumptions need not hold, even though the gradient-difference diagnostic remains computable.

For plain SGD, $u_{t-1}=-\eta g_{t-1}$ gives
\begin{equation}
 \widehat\kappa_t^{\mathrm{step}}
 =\frac{1}{\eta}\left(1-\frac{g_{t-1}^{\top}g_t}{\|g_{t-1}\|^2}\right)
 =\widehat\kappa_t.
\end{equation}
Its sampling term is $-\varepsilon_t/\eta$, where $\varepsilon_t$ is defined in the main text. In the same specialization the starting-point error bound becomes $L_H\eta\|g_{t-1}\|/2$. For a different optimizer, the raw-gradient combination $\eta\widehat\kappa_t$ and the step-direction curvature are generally different: one projects the gradient change onto the preceding gradient, the other onto the applied increment. The Adam overview retains the former measured quantity and labels it $\eta\widehat\kappa_t$; it is not relabeled as an empirical measurement of the latter.

A zero applied step makes $\widehat\kappa_t^{\mathrm{step}}$ undefined. Separately, $\eta\widehat\kappa_t$ is undefined when $g_{t-1}=0$; if only $g_t=0$, the inner-product expression gives $\eta\widehat\kappa_t=1$ although the gradient cosine is undefined. We use $r_t=\|g_t\|/\|g_{t-1}\|$ as the norm ratio in the supplementary timing tables.

\subsection{Sampling identities and their limits}
Fresh conditional unbiasedness gives $\E[g_{t-1}^\top\xi_t]=0$. At the preceding state, $g_{t-1}=f_{t-1}+\xi_{t-1}$ and $\E[f_{t-1}^\top\xi_{t-1}]=0$, hence
\begin{equation}
 \E[g_{t-1}^\top(\xi_t-\xi_{t-1})]=-\E\|\xi_{t-1}\|^2.
\end{equation}
This reasoning uses conditional expectations before and after the previous sampling decision; it does not assert independence between the new state and the previous noise. For independent uniform batches of $b$ distinct examples from $N$, standard finite-population sampling gives $\operatorname{Cov}(\xi\mid\theta)=(N-b)\Sigma/[b(N-1)]$, where $\Sigma$ uses denominator $N$. Taking its trace yields the variance identity.

At a frozen state, two independent batch gradients satisfy $\E[g^\top g']=\|f\|^2$ and $\E\|g\|^2=\|f\|^2+V_b$. Replacing the product of random norms by $\|f\|^2+V_b$ gives the approximate cosine formula. Neither division of expectations nor replacement of an evolving trajectory by a frozen one is exact. The fixed-state experiment tests these distinctions directly.

\subsection{Finite-step loss and an explicit residual}
Let $x=\theta_{t-1}$, $u=-\eta g_{t-1}$ and $\phi(a)=L(x+a u)$. Then $L_t-L_{t-1}=\int_0^1\phi'(a)\,da$. Integration by parts gives the exact endpoint-trapezoid residual
\begin{equation}
 \mathcal E_t:=L_t-L_{t-1}-\tfrac12(f_{t-1}+f_t)^\top u
 =\int_0^1(\tfrac12-a)u^\top H(x+a u)u\,da.
 \label{eq:t-trapezoid}
\end{equation}
Substituting $f_t=g_t-\xi_t$ yields the main sampled-loss identity. If $H$ is constant, the integral vanishes and full-batch GD gives $\Delta L_t=\eta\|g_{t-1}\|^2(\eta\widehat\kappa_t-2)/2$ exactly. If the Hessian is $L_H$-Lipschitz, subtract its value at the midpoint inside~\eqref{eq:t-trapezoid}; the constant term integrates to zero. Consequently,
\begin{equation}
 |\mathcal E_t|\leq L_H\|u\|^3\int_0^1|a-\tfrac12|^2\,da
 =L_H\|u\|^3/12.
\end{equation}
A small curvature discrepancy does not by itself imply a small error relative to a near-zero loss change. This motivates the separate empirical sign audit. The strict quadratic condition $\eta\widehat\kappa_t>2$ implies $\rho_t<0$ and $r_t>1$, but those two signs alone do not imply it. Current zero gradients and previous zero gradients retain the undefined-case rules in the main text.

\subsection{Adaptive updates}
For full-objective Adam, use $u_t=-\eta_t d_t$, with the parameter-state convention
\begin{align}
 m_t&=\beta_1m_{t-1}+(1-\beta_1)h_t,&
 v_t&=\beta_2v_{t-1}+(1-\beta_2)h_t^{\odot2},\notag\\
 d_t&=\frac{m_t/(1-\beta_1^{t+1})}{\sqrt{v_t/(1-\beta_2^{t+1})}+\epsilon},&
 \theta_{t+1}&=\theta_t-\eta_t d_t.
\end{align}
Moments are initialized to zero at index $-1$. Here $h_t=g_t$ for native Adam and $h_t=\min(1,\gamma/\|g_t\|)g_t$ for global clipping, with factor one when $g_t=0$. The raw diagnostic always uses the unclipped $g_t$. Late $\beta_2$ interventions preserve moments and the native counter; their bias correction uses the new coefficient and that existing counter, rather than a product of historical coefficients. Caps subsequently transform the stored parameter proposal while retaining these advanced moments. The step-direction estimate uses the actual increment after any intervention. The raw-gradient quantity $\eta\widehat\kappa_t$ continues to describe gradient reversal and growth, with the distinct interpretation explained above.
